% Options for packages loaded elsewhere
\PassOptionsToPackage{unicode}{hyperref}
\PassOptionsToPackage{hyphens}{url}
%
\documentclass[
  ignorenonframetext,
  aspectratio=169,
]{beamer}
\usepackage{pgfpages}
\setbeamertemplate{caption}[numbered]
\setbeamertemplate{caption label separator}{: }
\setbeamercolor{caption name}{fg=normal text.fg}
\beamertemplatenavigationsymbolshorizontal
% Prevent slide breaks in the middle of a paragraph
\widowpenalties 1 10000
\raggedbottom
\setbeamertemplate{part page}{
  \centering
  \begin{beamercolorbox}[sep=16pt,center]{part title}
    \usebeamerfont{part title}\insertpart\par
  \end{beamercolorbox}
}
\setbeamertemplate{section page}{
  \centering
  \begin{beamercolorbox}[sep=12pt,center]{part title}
    \usebeamerfont{section title}\insertsection\par
  \end{beamercolorbox}
}
\setbeamertemplate{subsection page}{
  \centering
  \begin{beamercolorbox}[sep=8pt,center]{part title}
    \usebeamerfont{subsection title}\insertsubsection\par
  \end{beamercolorbox}
}
\AtBeginPart{
  \frame{\partpage}
}
\AtBeginSection{
  \ifbibliography
  \else
    \frame{\sectionpage}
  \fi
}
\AtBeginSubsection{
  \frame{\subsectionpage}
}

\usepackage{amsmath,amssymb}
\usepackage{iftex}
\ifPDFTeX
  \usepackage[T1]{fontenc}
  \usepackage[utf8]{inputenc}
  \usepackage{textcomp} % provide euro and other symbols
\else % if luatex or xetex
  \usepackage{unicode-math}
  \defaultfontfeatures{Scale=MatchLowercase}
  \defaultfontfeatures[\rmfamily]{Ligatures=TeX,Scale=1}
\fi
\usepackage{lmodern}
\ifPDFTeX\else  
    % xetex/luatex font selection
\fi
% Use upquote if available, for straight quotes in verbatim environments
\IfFileExists{upquote.sty}{\usepackage{upquote}}{}
\IfFileExists{microtype.sty}{% use microtype if available
  \usepackage[]{microtype}
  \UseMicrotypeSet[protrusion]{basicmath} % disable protrusion for tt fonts
}{}
\usepackage{xcolor}
\newif\ifbibliography
\setlength{\emergencystretch}{3em} % prevent overfull lines
\setcounter{secnumdepth}{-\maxdimen} % remove section numbering

\usepackage{color}
\usepackage{fancyvrb}
\newcommand{\VerbBar}{|}
\newcommand{\VERB}{\Verb[commandchars=\\\{\}]}
\DefineVerbatimEnvironment{Highlighting}{Verbatim}{commandchars=\\\{\}}
% Add ',fontsize=\small' for more characters per line
\usepackage{framed}
\definecolor{shadecolor}{RGB}{241,243,245}
\newenvironment{Shaded}{\begin{snugshade}}{\end{snugshade}}
\newcommand{\AlertTok}[1]{\textcolor[rgb]{0.68,0.00,0.00}{#1}}
\newcommand{\AnnotationTok}[1]{\textcolor[rgb]{0.37,0.37,0.37}{#1}}
\newcommand{\AttributeTok}[1]{\textcolor[rgb]{0.40,0.45,0.13}{#1}}
\newcommand{\BaseNTok}[1]{\textcolor[rgb]{0.68,0.00,0.00}{#1}}
\newcommand{\BuiltInTok}[1]{\textcolor[rgb]{0.00,0.23,0.31}{#1}}
\newcommand{\CharTok}[1]{\textcolor[rgb]{0.13,0.47,0.30}{#1}}
\newcommand{\CommentTok}[1]{\textcolor[rgb]{0.37,0.37,0.37}{#1}}
\newcommand{\CommentVarTok}[1]{\textcolor[rgb]{0.37,0.37,0.37}{\textit{#1}}}
\newcommand{\ConstantTok}[1]{\textcolor[rgb]{0.56,0.35,0.01}{#1}}
\newcommand{\ControlFlowTok}[1]{\textcolor[rgb]{0.00,0.23,0.31}{\textbf{#1}}}
\newcommand{\DataTypeTok}[1]{\textcolor[rgb]{0.68,0.00,0.00}{#1}}
\newcommand{\DecValTok}[1]{\textcolor[rgb]{0.68,0.00,0.00}{#1}}
\newcommand{\DocumentationTok}[1]{\textcolor[rgb]{0.37,0.37,0.37}{\textit{#1}}}
\newcommand{\ErrorTok}[1]{\textcolor[rgb]{0.68,0.00,0.00}{#1}}
\newcommand{\ExtensionTok}[1]{\textcolor[rgb]{0.00,0.23,0.31}{#1}}
\newcommand{\FloatTok}[1]{\textcolor[rgb]{0.68,0.00,0.00}{#1}}
\newcommand{\FunctionTok}[1]{\textcolor[rgb]{0.28,0.35,0.67}{#1}}
\newcommand{\ImportTok}[1]{\textcolor[rgb]{0.00,0.46,0.62}{#1}}
\newcommand{\InformationTok}[1]{\textcolor[rgb]{0.37,0.37,0.37}{#1}}
\newcommand{\KeywordTok}[1]{\textcolor[rgb]{0.00,0.23,0.31}{\textbf{#1}}}
\newcommand{\NormalTok}[1]{\textcolor[rgb]{0.00,0.23,0.31}{#1}}
\newcommand{\OperatorTok}[1]{\textcolor[rgb]{0.37,0.37,0.37}{#1}}
\newcommand{\OtherTok}[1]{\textcolor[rgb]{0.00,0.23,0.31}{#1}}
\newcommand{\PreprocessorTok}[1]{\textcolor[rgb]{0.68,0.00,0.00}{#1}}
\newcommand{\RegionMarkerTok}[1]{\textcolor[rgb]{0.00,0.23,0.31}{#1}}
\newcommand{\SpecialCharTok}[1]{\textcolor[rgb]{0.37,0.37,0.37}{#1}}
\newcommand{\SpecialStringTok}[1]{\textcolor[rgb]{0.13,0.47,0.30}{#1}}
\newcommand{\StringTok}[1]{\textcolor[rgb]{0.13,0.47,0.30}{#1}}
\newcommand{\VariableTok}[1]{\textcolor[rgb]{0.07,0.07,0.07}{#1}}
\newcommand{\VerbatimStringTok}[1]{\textcolor[rgb]{0.13,0.47,0.30}{#1}}
\newcommand{\WarningTok}[1]{\textcolor[rgb]{0.37,0.37,0.37}{\textit{#1}}}

\providecommand{\tightlist}{%
  \setlength{\itemsep}{0pt}\setlength{\parskip}{0pt}}\usepackage{longtable,booktabs,array}
\usepackage{calc} % for calculating minipage widths
\usepackage{caption}
% Make caption package work with longtable
\makeatletter
\def\fnum@table{\tablename~\thetable}
\makeatother
\usepackage{graphicx}
\makeatletter
\def\maxwidth{\ifdim\Gin@nat@width>\linewidth\linewidth\else\Gin@nat@width\fi}
\def\maxheight{\ifdim\Gin@nat@height>\textheight\textheight\else\Gin@nat@height\fi}
\makeatother
% Scale images if necessary, so that they will not overflow the page
% margins by default, and it is still possible to overwrite the defaults
% using explicit options in \includegraphics[width, height, ...]{}
\setkeys{Gin}{width=\maxwidth,height=\maxheight,keepaspectratio}
% Set default figure placement to htbp
\makeatletter
\def\fps@figure{htbp}
\makeatother

%%% Load theme
\IfFileExists{beamerthemegotham.sty}{
  %%%% Full TeXlive
  \usetheme{gotham}
  \gothamset{
    numbering=totalframenumber,
    parttocframe default=off,
    sectiontocframe default=off,
    subsectiontocframe default=off,
    tocframe template=gotham bullet,
    % progressbar position=frametitle,
    progressbar position=foot,
  }
  \renewcommand{\gothamInstituteLogoSquare}[1][4ex]{%
    \includegraphics[height=##1]{_resources/image/logo_rudn-inverted.png}
  }
  % \renewcommand{\gothamtitlepagelogo}[1][4ex]{%
  %   \includegraphics[height=##1]{_resources/image/logo_rudn.png}
  % }
}{%
  %%%% TinyTeX
  \usepackage{libertine}
  %%%% https://deic.uab.cat/~iblanes/beamer_gallery/
  \usetheme{Madrid}
}
\makeatletter
\@ifpackageloaded{caption}{}{\usepackage{caption}}
\AtBeginDocument{%
\ifdefined\contentsname
  \renewcommand*\contentsname{Содержание}
\else
  \newcommand\contentsname{Содержание}
\fi
\ifdefined\listfigurename
  \renewcommand*\listfigurename{Список иллюстраций}
\else
  \newcommand\listfigurename{Список иллюстраций}
\fi
\ifdefined\listtablename
  \renewcommand*\listtablename{Список таблиц}
\else
  \newcommand\listtablename{Список таблиц}
\fi
\ifdefined\figurename
  \renewcommand*\figurename{Рисунок}
\else
  \newcommand\figurename{Рисунок}
\fi
\ifdefined\tablename
  \renewcommand*\tablename{Таблица}
\else
  \newcommand\tablename{Таблица}
\fi
}
\@ifpackageloaded{float}{}{\usepackage{float}}
\floatstyle{ruled}
\@ifundefined{c@chapter}{\newfloat{codelisting}{h}{lop}}{\newfloat{codelisting}{h}{lop}[chapter]}
\floatname{codelisting}{Список}
\newcommand*\listoflistings{\listof{codelisting}{Листинги}}
\makeatother
\makeatletter
\makeatother
\makeatletter
\@ifpackageloaded{caption}{}{\usepackage{caption}}
\@ifpackageloaded{subcaption}{}{\usepackage{subcaption}}
\makeatother

\ifLuaTeX
\usepackage[bidi=basic]{babel}
\else
\usepackage[bidi=default]{babel}
\fi
\babelprovide[main,import]{russian}
\babelprovide[import]{english}
% get rid of language-specific shorthands (see #6817):
\let\LanguageShortHands\languageshorthands
\def\languageshorthands#1{}
\ifLuaTeX
  \usepackage{selnolig}  % disable illegal ligatures
\fi
\usepackage{csquotes}
\usepackage{bookmark}

\IfFileExists{xurl.sty}{\usepackage{xurl}}{} % add URL line breaks if available
\urlstyle{same} % disable monospaced font for URLs
\hypersetup{
  pdftitle={Научная презентация},
  pdfauthor={Дмитрий Сергеевич Кулябов},
  pdflang={ru-RU},
  hidelinks,
  pdfcreator={LaTeX via pandoc}}


\title{Научная презентация}
\subtitle{Структура и оформление}
\author{Дмитрий Сергеевич Кулябов}
\date{2026-10-13}

\begin{document}
\frame{\titlepage}

\renewcommand*\contentsname{Содержание}
\begin{frame}[allowframebreaks]
  \frametitle{Содержание}
  \tableofcontents[hideallsubsections]
\end{frame}

\section{Информация}\label{ux438ux43dux444ux43eux440ux43cux430ux446ux438ux44f}

\begin{frame}{Докладчик}
\phantomsection\label{ux434ux43eux43aux43bux430ux434ux447ux438ux43a}
\begin{columns}[c]
\begin{column}{0.7\textwidth}
\begin{itemize}
\tightlist
\item
  Кулябов Дмитрий Сергеевич
\item
  д.ф.-м.н., профессор
\item
  профессор кафедры теории вероятностей и кибербезопасности
\item
  Российский университет дружбы народов им. П. Лумумбы
\item
  \href{mailto:kulyabov-ds@rudn.ru}{\nolinkurl{kulyabov-ds@rudn.ru}}
\item
  \url{https://yamadharma.github.io/ru/}
\end{itemize}
\end{column}

\begin{column}{0.3\textwidth}
\includegraphics{./image/kulyabov.jpg}
\end{column}
\end{columns}

\note{\begin{itemize}
\tightlist
\item
  Поприветствовать аудиторию.
\item
  Представиться кратко, даже если вас объявили.
\end{itemize}}
\end{frame}

\section{Вводная
часть}\label{ux432ux432ux43eux434ux43dux430ux44f-ux447ux430ux441ux442ux44c}

\begin{frame}{Актуальность}
\phantomsection\label{ux430ux43aux442ux443ux430ux43bux44cux43dux43eux441ux442ux44c}
\begin{itemize}
\tightlist
\item
  Важно донести результаты своих исследований до окружающих
\item
  Научная презентация --- рабочий инструмент исследователя
\item
  Необходимо создавать презентацию быстро
\item
  Желательна минимизация усилий для создания презентации
\end{itemize}

\note{\begin{itemize}
\tightlist
\item
  Почему именно \emph{структура} презентации важна.
\item
  Время: не более 1 минуты на этот слайд.
\end{itemize}}
\end{frame}

\begin{frame}{Объект и предмет исследования}
\phantomsection\label{ux43eux431ux44aux435ux43aux442-ux438-ux43fux440ux435ux434ux43cux435ux442-ux438ux441ux441ux43bux435ux434ux43eux432ux430ux43dux438ux44f}
\begin{itemize}
\tightlist
\item
  Презентация как текст
\item
  Программное обеспечение для создания презентаций
\item
  Входные и выходные форматы презентаций
\end{itemize}

\note{\begin{itemize}
\tightlist
\item
  Объект --- более широкая область (презентации в науке).
\item
  Предмет --- конкретный аспект (автоматизация создания через Markdown).
\item
  Рассматриваем презентацию не как на набор слайдов в PowerPoint, а как
  на код.
\end{itemize}}
\end{frame}

\begin{frame}{Цели и задачи}
\phantomsection\label{ux446ux435ux43bux438-ux438-ux437ux430ux434ux430ux447ux438}
\begin{itemize}
\tightlist
\item
  Цель:

  \begin{itemize}
  \tightlist
  \item
    Разработать эффективный и воспроизводимый шаблон для научных
    презентаций.
  \end{itemize}
\item
  Задачи:

  \begin{itemize}
  \tightlist
  \item
    Создать шаблон презентации в Markdown.
  \item
    Описать алгоритм сборки презентации в форматы PDF и HTML.
  \item
    Сформулировать рекомендации по наполнению слайдов.
  \end{itemize}
\end{itemize}

\note{\begin{itemize}
\tightlist
\item
  Обратить внимание, задач ровно три.
\item
  Это удобно для восприятия.
\end{itemize}}
\end{frame}

\begin{frame}[fragile]{Материалы и методы}
\phantomsection\label{ux43cux430ux442ux435ux440ux438ux430ux43bux44b-ux438-ux43cux435ux442ux43eux434ux44b}
\begin{itemize}
\tightlist
\item
  Процессор \texttt{pandoc} (как часть \texttt{quarto}) для входного
  формата Markdown
\item
  Результирующие форматы

  \begin{itemize}
  \tightlist
  \item
    \texttt{pdf}
  \item
    \texttt{html}
  \end{itemize}
\item
  Автоматизация процесса создания: \texttt{Makefile}
\end{itemize}

\note{\begin{itemize}
\tightlist
\item
  Предлагается использовать наиболее популярные форматы.
\item
  \texttt{Makefile} позволяет одной командой собрать все форматы.
\end{itemize}}
\end{frame}

\section{Создание
презентации}\label{ux441ux43eux437ux434ux430ux43dux438ux435-ux43fux440ux435ux437ux435ux43dux442ux430ux446ux438ux438}

\begin{frame}[fragile]{Пакет \texttt{quarto}}
\phantomsection\label{ux43fux430ux43aux435ux442-quarto}
\begin{itemize}
\tightlist
\item
  Quarto --- современная система для создания научной, технической и
  прочей документации.
\item
  Сайт: \url{https://quarto.org/}.
\item
  Репозиторий: \url{https://github.com/quarto-dev/quarto-cli}.
\item
  Входной язык: Markdown.
\item
  Выходные форматы: html, pdf, epub, docx, презентации в формате
  reveal.js.
\item
  Интеграция с языками программирования: R, Python, Julia.
\end{itemize}

\note{\begin{itemize}
\tightlist
\item
  Пишем один раз, а получаем слайды в HTML и PDF.
\item
  Демонстрация: Можно показать на втором экране консоль с командой
  \texttt{make}.
\end{itemize}}
\end{frame}

\begin{frame}[fragile]{Процессор \texttt{pandoc}}
\phantomsection\label{ux43fux440ux43eux446ux435ux441ux441ux43eux440-pandoc}
\begin{itemize}
\tightlist
\item
  Один из движков для \texttt{quato}.
\item
  Pandoc: преобразователь текстовых файлов.
\item
  Сайт: \url{https://pandoc.org/}.
\item
  Репозиторий: \url{https://github.com/jgm/pandoc}.
\end{itemize}

\note{\begin{itemize}
\tightlist
\item
  Quarto = Pandoc на стероидах.
\end{itemize}}
\end{frame}

\begin{frame}[fragile]{Формат \texttt{pdf}}
\phantomsection\label{ux444ux43eux440ux43cux430ux442-pdf}
\begin{itemize}
\tightlist
\item
  Использование LaTeX
\item
  Пакет для презентации: \href{https://ctan.org/pkg/beamer}{beamer}
\item
  Тема оформления: \texttt{gotham} как развитие \texttt{metropolis}
\end{itemize}

\note{\begin{itemize}
\tightlist
\item
  Преимущество PDF:

  \begin{itemize}
  \tightlist
  \item
    универсальный формат;
  \item
    не зависит от версии PowerPoint;
  \item
    шрифты встроены;
  \item
    формулы выглядят идеально.
  \end{itemize}
\item
  Тема \texttt{gotham} как развитие \texttt{metropolis}.

  \begin{itemize}
  \tightlist
  \item
    Более современный вид презентации.
  \end{itemize}
\end{itemize}}
\end{frame}

\begin{frame}[fragile]{Код для формата \texttt{pdf}}
\phantomsection\label{ux43aux43eux434-ux434ux43bux44f-ux444ux43eux440ux43cux430ux442ux430-pdf}
\begin{Shaded}
\begin{Highlighting}[]
\FunctionTok{slide\_level}\KeywordTok{:}\AttributeTok{ }\DecValTok{2}
\FunctionTok{aspectratio}\KeywordTok{:}\AttributeTok{ }\DecValTok{169}
\FunctionTok{section{-}titles}\KeywordTok{:}\AttributeTok{ }\CharTok{true}
\end{Highlighting}
\end{Shaded}

\note{\begin{itemize}
\tightlist
\item
  Уровень слайда (2 или 3).
\item
  Соотношение сторон (16x9).
\end{itemize}}
\end{frame}

\begin{frame}[fragile]{Формат \texttt{html}}
\phantomsection\label{ux444ux43eux440ux43cux430ux442-html}
\begin{itemize}
\tightlist
\item
  Используется фреймворк \href{https://revealjs.com/}{reveal.js}
\item
  Презентация работает в браузере, не требует интернета (все ресурсы
  упакованы)
\item
  Используется \href{https://revealjs.com/themes/}{тема} \texttt{beige}
  (светлая, подходит для хорошо освещенных аудиторий)
\end{itemize}

\note{\begin{itemize}
\tightlist
\item
  Интерактивность.

  \begin{itemize}
  \tightlist
  \item
    В HTML-версии можно использовать фрагменты (появление по клику), что
    удобно для управления вниманием.
  \end{itemize}
\item
  Портативность.

  \begin{itemize}
  \tightlist
  \item
    Папку с презентацией можно просто скопировать на флешку и открыть на
    любом компьютере.
  \end{itemize}
\end{itemize}}
\end{frame}

\section{Результаты}\label{ux440ux435ux437ux443ux43bux44cux442ux430ux442ux44b}

\begin{frame}[fragile]{Получающиеся форматы}
\phantomsection\label{ux43fux43eux43bux443ux447ux430ux44eux449ux438ux435ux441ux44f-ux444ux43eux440ux43cux430ux442ux44b}
\begin{itemize}
\tightlist
\item
  Полученный \texttt{pdf}-файл можно демонстрировать в любой программе
  просмотра \texttt{pdf}
\item
  Полученный \texttt{html}-файл содержит в себе все ресурсы:
  изображения, css, скрипты
\end{itemize}

\note{\begin{itemize}
\tightlist
\item
  Воспроизводимость: Исходный Markdown легко хранить в системе контроля
  версий.
\item
  Итог раздела: Мы автоматизировали рутину и можем сосредоточиться на
  содержании.
\item
  Вопрос на подумать: Сколько времени вы тратите на расстановку картинок
  в PowerPoint вручную?
\end{itemize}}
\end{frame}

\section{Элементы
презентации}\label{ux44dux43bux435ux43cux435ux43dux442ux44b-ux43fux440ux435ux437ux435ux43dux442ux430ux446ux438ux438}

\begin{frame}{Актуальность}
\phantomsection\label{ux430ux43aux442ux443ux430ux43bux44cux43dux43eux441ux442ux44c-1}
\begin{itemize}
\tightlist
\item
  Даёт понять, о чём пойдёт речь
\item
  Следует широко и кратко описать проблему
\item
  Мотивировать свое исследование
\item
  Сформулировать цели и задачи
\item
  Возможна формулировка ожидаемых результатов
\end{itemize}

\note{\begin{itemize}
\tightlist
\item
  Широкий контекст.

  \begin{itemize}
  \tightlist
  \item
    Опишите проблему, понятную не только узким специалистам.
  \end{itemize}
\item
  Мотивация

  \begin{itemize}
  \tightlist
  \item
    Почему ваше исследование важно именно сейчас.
  \end{itemize}
\item
  Цель и ожидаемый результат

  \begin{itemize}
  \tightlist
  \item
    Сформулируйте четко, что слушатель получит в конце.
  \end{itemize}
\item
  Начните с вопроса к залу или яркого факта.

  \begin{itemize}
  \tightlist
  \item
    Например: \enquote{Знаете ли вы, что среднее время удержания
    внимания взрослого человека --- 8 секунд?}
  \end{itemize}
\item
  Избегайте

  \begin{itemize}
  \tightlist
  \item
    Слишком общих фраз (например, это важно для науки). Конкретика!
  \end{itemize}
\end{itemize}}
\end{frame}

\begin{frame}{Цели и задачи}
\phantomsection\label{ux446ux435ux43bux438-ux438-ux437ux430ux434ux430ux447ux438-1}
\begin{itemize}
\tightlist
\item
  Правило 1--2

  \begin{itemize}
  \tightlist
  \item
    Не ставьте более двух целей на одну презентацию.
  \end{itemize}
\item
  SMART

  \begin{itemize}
  \tightlist
  \item
    Specific, Measurable, Achievable, Relevant, Time-bound.
  \end{itemize}
\item
  Цель --- что хотим достичь.
\item
  Задачи --- конкретные шаги.
\end{itemize}

\note{\begin{itemize}
\tightlist
\item
  Пример плохо

  \begin{itemize}
  \tightlist
  \item
    \enquote{Изучить свойства графена.} (слишком широко).
  \end{itemize}
\item
  Пример хорошо

  \begin{itemize}
  \tightlist
  \item
    \enquote{Показать, что предложенный метод синтеза повышает прочность
    композита на 15\%.}
  \end{itemize}
\end{itemize}}
\end{frame}

\begin{frame}{Материалы и методы}
\phantomsection\label{ux43cux430ux442ux435ux440ux438ux430ux43bux44b-ux438-ux43cux435ux442ux43eux434ux44b-1}
\begin{itemize}
\tightlist
\item
  Представляйте данные качественно
\item
  Количественно, только если крайне необходимо
\item
  Качественно \textgreater{} Количественно

  \begin{itemize}
  \tightlist
  \item
    Показывайте схемы установок, блок-схемы алгоритмов.
  \end{itemize}
\item
  Детализация

  \begin{itemize}
  \tightlist
  \item
    Излишние детали не нужны
  \item
    На слайде --- принцип.
  \item
    В речи --- детали (или на раздаточном материале).
  \end{itemize}
\item
  Таблицы

  \begin{itemize}
  \tightlist
  \item
    Не более 4 столбцов и 5 строк на слайде.
  \end{itemize}
\end{itemize}

\note{\begin{itemize}
\tightlist
\item
  Если много данных

  \begin{itemize}
  \tightlist
  \item
    Вынесите подробную таблицу в раздаточный материал (handout)
  \item
    Или скажите: \enquote{Подробные параметры приведены в
    статье/диссертации}.
  \end{itemize}
\item
  Анимация

  \begin{itemize}
  \tightlist
  \item
    Сложную схему можно показывать по частям, используя фрагменты.
  \end{itemize}
\end{itemize}}
\end{frame}

\begin{frame}{Содержание исследования}
\phantomsection\label{ux441ux43eux434ux435ux440ux436ux430ux43dux438ux435-ux438ux441ux441ux43bux435ux434ux43eux432ux430ux43dux438ux44f}
\begin{itemize}
\tightlist
\item
  Предлагаемое решение задач исследования с обоснованием.
\item
  Основные этапы работы.
\item
  Логика

  \begin{itemize}
  \tightlist
  \item
    Проблема → Гипотеза → Эксперимент → Анализ → Вывод.
  \end{itemize}
\item
  Один слайд --- один этап

  \begin{itemize}
  \tightlist
  \item
    Не пытайтесь уместить всё на одном экране.
  \end{itemize}
\item
  Необходимы логические связки между слайдами

  \begin{itemize}
  \tightlist
  \item
    Используйте фразы-переходы между слайдами.
  \end{itemize}
\end{itemize}

\note{\begin{itemize}
\tightlist
\item
  Навигация

  \begin{itemize}
  \tightlist
  \item
    Полезно иметь в углу слайда миниатюру структуры доклада с подсветкой
    текущего раздела (в beamer это делается автоматически).
  \item
    Страницы в формате страница/всего страниц.
  \end{itemize}
\item
  Речь

  \begin{itemize}
  \tightlist
  \item
    \enquote{Итак, мы увидели проблему. Теперь посмотрим, как мы
    предлагаем её решить\ldots{}}
  \item
    Это помогает аудитории не потеряться.
  \end{itemize}
\end{itemize}}
\end{frame}

\begin{frame}{Результаты}
\phantomsection\label{ux440ux435ux437ux443ux43bux44cux442ux430ux442ux44b-1}
\begin{itemize}
\tightlist
\item
  Не нужны все результаты
\item
  Принцип Парето

  \begin{itemize}
  \tightlist
  \item
    Покажите 20\% результатов, которые дают 80\% понимания.
  \end{itemize}
\item
  Визуализация

  \begin{itemize}
  \tightlist
  \item
    График или диаграмма всегда лучше столбца цифр.
  \end{itemize}
\item
  Интерпретация

  \begin{itemize}
  \tightlist
  \item
    Подпишите оси, выделите ключевую точку на графике.
  \end{itemize}
\item
  Необходимо показать понимание материала
\end{itemize}

\note{\begin{itemize}
\tightlist
\item
  Золотое правило

  \begin{itemize}
  \tightlist
  \item
    На слайде с результатами должна быть одна главная мысль.
  \item
    Если мыслей две --- делайте два слайда.
  \end{itemize}
\item
  Цвет

  \begin{itemize}
  \tightlist
  \item
    Используйте контрастный цвет, чтобы выделить главную кривую на
    графике.
  \end{itemize}
\end{itemize}}
\end{frame}

\begin{frame}{Итоговый слайд}
\phantomsection\label{ux438ux442ux43eux433ux43eux432ux44bux439-ux441ux43bux430ux439ux434}
\begin{itemize}
\tightlist
\item
  Запоминается последняя фраза. © Штирлиц
\item
  Главное сообщение, которое вы хотите донести до слушателей

  \begin{itemize}
  \tightlist
  \item
    1--2 предложения, которые слушатель должен запомнить.
  \end{itemize}
\item
  Избегайте использовать последний слайд вида

  \begin{itemize}
  \tightlist
  \item
    \emph{Спасибо за внимание}
  \item
    \emph{Вопросы?}
  \end{itemize}
\end{itemize}

\note{\begin{itemize}
\tightlist
\item
  Пример сильной концовки

  \begin{itemize}
  \tightlist
  \item
    \enquote{Таким образом, использование Markdown сокращает время
    подготовки презентации в 3 раза, позволяя ученому заниматься наукой,
    а не версткой.}
  \end{itemize}
\item
  После этого спокойно скажите

  \begin{itemize}
  \tightlist
  \item
    \enquote{Я готов ответить на ваши вопросы.}
  \end{itemize}
\end{itemize}}
\end{frame}

\section{Рекомендую}\label{ux440ux435ux43aux43eux43cux435ux43dux434ux443ux44e}

\begin{frame}{Принцип 10/20/30 (Гай Кавасаки)}
\phantomsection\label{ux43fux440ux438ux43dux446ux438ux43f-102030-ux433ux430ux439-ux43aux430ux432ux430ux441ux430ux43aux438}
\begin{itemize}
\tightlist
\item
  10 слайдов

  \begin{itemize}
  \tightlist
  \item
    Оптимально для 20-минутного выступления.
  \end{itemize}
\item
  20 минут

  \begin{itemize}
  \tightlist
  \item
    Остальное время --- обсуждение.
  \end{itemize}
\item
  30 кегль

  \begin{itemize}
  \tightlist
  \item
    Минимальный размер шрифта для читаемости с задних рядов.
  \end{itemize}
\end{itemize}

\note{\begin{itemize}
\tightlist
\item
  Адаптация

  \begin{itemize}
  \tightlist
  \item
    Если регламент 15 минут --- 7-8 слайдов.
  \item
    Если 45 минут (лекция) --- 20-25 слайдов.
  \end{itemize}
\item
  Проверка шрифта

  \begin{itemize}
  \tightlist
  \item
    Распечатайте слайд на листе А4 и положите на пол.
  \item
    Если с высоты роста не видно --- шрифт мелкий.
  \end{itemize}
\end{itemize}}
\end{frame}

\begin{frame}{Связь слайдов (нарратив)}
\phantomsection\label{ux441ux432ux44fux437ux44c-ux441ux43bux430ux439ux434ux43eux432-ux43dux430ux440ux440ux430ux442ux438ux432}
\begin{itemize}[<+->]
\tightlist
\item
  Один слайд --- одна мысль
\item
  Не ссылайтесь назад

  \begin{itemize}[<+->]
  \tightlist
  \item
    Нельзя ссылаться на объекты, находящиеся на предыдущих слайдах
    (например, на формулы)
  \item
    Слушатель не должен помнить формулу с предыдущего слайда.
  \item
    Либо повторите ее, либо вынесите на доску.
  \end{itemize}
\item
  Каждый слайд должен иметь заголовок.
\item
  Storytelling

  \begin{itemize}[<+->]
  \tightlist
  \item
    У презентации должна быть завязка (проблема), развитие
    (исследование) и развязка (решение).
  \end{itemize}
\end{itemize}

\note{\begin{itemize}
\tightlist
\item
  Техника Понятно / Непонятно

  \begin{itemize}
  \tightlist
  \item
    После каждого логического блока мысленно спросите: \enquote{Всем ли
    понятно?}.
  \item
    Если нет --- сделайте паузу или приведите аналогию.
  \end{itemize}
\item
  Паузы

  \begin{itemize}
  \tightlist
  \item
    Паузы важнее слов.
  \item
    После ключевого вывода сделайте паузу в 3-5 секунд.
  \end{itemize}
\end{itemize}}
\end{frame}

\begin{frame}{Количество сущностей (когнитивная нагрузка)}
\phantomsection\label{ux43aux43eux43bux438ux447ux435ux441ux442ux432ux43e-ux441ux443ux449ux43dux43eux441ux442ux435ux439-ux43aux43eux433ux43dux438ux442ux438ux432ux43dux430ux44f-ux43dux430ux433ux440ux443ux437ux43aux430}
\begin{itemize}[<+->]
\tightlist
\item
  Правило \(7 \pm 2\)

  \begin{itemize}[<+->]
  \tightlist
  \item
    Человек удерживает в оперативной памяти 5--9 объектов.
  \end{itemize}
\item
  Группировка

  \begin{itemize}[<+->]
  \tightlist
  \item
    Объединяйте элементы в визуальные блоки (рамки, отступы).
  \end{itemize}
\item
  Правило 5

  \begin{itemize}[<+->]
  \tightlist
  \item
    Стремитесь, чтобы на слайде было не более 5 смысловых блоков.
  \item
    Можно группировать элементы так, чтобы визуально было не более 5
    групп
  \end{itemize}
\end{itemize}

\note{\begin{itemize}
\tightlist
\item
  Пример

  \begin{itemize}
  \tightlist
  \item
    Список из 10 пунктов разбейте на две колонки по 5.
  \end{itemize}
\item
  Фрагменты

  \begin{itemize}
  \tightlist
  \item
    Показывайте пункты списка последовательно, а не все сразу.
  \item
    Это управляет вниманием.
  \end{itemize}
\end{itemize}}
\end{frame}

\begin{frame}[fragile]{Общие рекомендации (дизайн и подача)}
\phantomsection\label{ux43eux431ux449ux438ux435-ux440ux435ux43aux43eux43cux435ux43dux434ux430ux446ux438ux438-ux434ux438ux437ux430ux439ux43d-ux438-ux43fux43eux434ux430ux447ux430}
\begin{itemize}[<+->]
\tightlist
\item
  Дополнение, а не дублирование

  \begin{itemize}[<+->]
  \tightlist
  \item
    Слайды --- это иллюстрация к речи, а не суфлёр для докладчика.
  \item
    На слайд выносится та информация, которая без зрительной опоры
    воспринимается хуже.
  \end{itemize}
\item
  Краткость

  \begin{itemize}[<+->]
  \tightlist
  \item
    Максимум 6-7 слов в строке, 6-7 строк на слайде.
  \end{itemize}
\item
  Графика

  \begin{itemize}[<+->]
  \tightlist
  \item
    Каждая картинка должна нести смысл.
  \item
    Избегайте клипарта и водяных знаков на всю площадь.
  \end{itemize}
\item
  Анимация

  \begin{itemize}[<+->]
  \tightlist
  \item
    Не злоупотребляйте анимацией и переходами
  \item
    Только для последовательного раскрытия информации.
  \item
    Избегайте вращений и вылетов.
  \end{itemize}
\end{itemize}

\note{\begin{itemize}
\tightlist
\item
  Второй экран (режим докладчика)

  \begin{itemize}
  \tightlist
  \item
    В \texttt{reveal.js} нажмите \texttt{S} (Speaker View).
  \item
    Вы увидите текущий слайд, следующий слайд, заметки и таймер.
  \item
    Это ваш главный инструмент на выступлении.
  \end{itemize}
\item
  Лазерная указка

  \begin{itemize}
  \tightlist
  \item
    Не водите ей хаотично по экрану.
  \item
    Обвели объект --- убрали руку.
  \end{itemize}
\end{itemize}}
\end{frame}

\begin{frame}{Иерархия представления данных}
\phantomsection\label{ux438ux435ux440ux430ux440ux445ux438ux44f-ux43fux440ux435ux434ux441ux442ux430ux432ux43bux435ux43dux438ux44f-ux434ux430ux43dux43dux44bux445}
\begin{itemize}[<+->]
\tightlist
\item
  Схема / Блок-схема

  \begin{itemize}[<+->]
  \tightlist
  \item
    Лучший выбор для объяснения процессов.
  \end{itemize}
\item
  График

  \begin{itemize}[<+->]
  \tightlist
  \item
    Для демонстрации зависимостей и трендов.
  \end{itemize}
\item
  Диаграмма

  \begin{itemize}[<+->]
  \tightlist
  \item
    Для сравнения величин.
  \end{itemize}
\item
  Таблица

  \begin{itemize}[<+->]
  \tightlist
  \item
    Только для точных значений (и то, минимум данных).
  \end{itemize}
\item
  Текст

  \begin{itemize}[<+->]
  \tightlist
  \item
    Крайний случай, когда без определения или формулы не обойтись.
  \end{itemize}
\end{itemize}

\note{\begin{itemize}
\tightlist
\item
  Цветовая слепота

  \begin{itemize}
  \tightlist
  \item
    Не используйте сочетания красный/зеленый для различения данных.
    Лучше синий/оранжевый.
  \end{itemize}
\item
  Подписи

  \begin{itemize}
  \tightlist
  \item
    Подпись к рисунку должна быть информативной, а не просто
    \enquote{Рис. 1}.
  \end{itemize}
\end{itemize}}
\end{frame}

\begin{frame}{Подготовка к выступлению}
\phantomsection\label{ux43fux43eux434ux433ux43eux442ux43eux432ux43aux430-ux43a-ux432ux44bux441ux442ux443ux43fux43bux435ux43dux438ux44e}
\begin{itemize}[<+->]
\tightlist
\item
  Репетиция

  \begin{itemize}[<+->]
  \tightlist
  \item
    Обязательно вслух и с таймером.
  \end{itemize}
\item
  План Б

  \begin{itemize}[<+->]
  \tightlist
  \item
    Сохраните презентацию в PDF и HTML.
  \item
    HTML может выручить, если нет средств просмотра pdf на чужом
    компьютере.
  \end{itemize}
\item
  Вопросы

  \begin{itemize}[<+->]
  \tightlist
  \item
    Заранее продумайте 2-3 слабых места в работе и подготовьте ответы на
    возможные вопросы.
  \end{itemize}
\end{itemize}

\note{\begin{itemize}
\tightlist
\item
  Дыхание

  \begin{itemize}
  \tightlist
  \item
    Перед выходом сделайте глубокий вдох.
  \item
    Первые 30 секунд --- самые волнительные.
  \end{itemize}
\item
  Контакт глаз

  \begin{itemize}
  \tightlist
  \item
    Смотрите на аудиторию, а не на экран ноутбука.
  \item
    Для этого и нужен второй экран с заметками.
  \end{itemize}
\end{itemize}}
\end{frame}

\section{Заключение}\label{ux437ux430ux43aux43bux44eux447ux435ux43dux438ux435}

\begin{frame}{Главный вывод}
\phantomsection\label{ux433ux43bux430ux432ux43dux44bux439-ux432ux44bux432ux43eux434}
Структурированная презентация на основе Markdown позволяет исследователю
сосредоточиться на науке, а не на борьбе с интерфейсом.

\note{\begin{itemize}
\tightlist
\item
  Финальный аккорд

  \begin{itemize}
  \tightlist
  \item
    Повторите главную мысль, поблагодарите за внимание (устно) и
    переходите к вопросам.
  \end{itemize}
\item
  Контакты

  \begin{itemize}
  \tightlist
  \item
    Оставьте слайд с вашим фото и email на время сессии вопросов.
  \end{itemize}
\end{itemize}}
\end{frame}




\end{document}
